\documentclass[11pt,reqno]{amsart}
\usepackage[T1]{fontenc}
\usepackage{lmodern}
\usepackage{microtype}
\usepackage{amsmath,amssymb,mathtools}
\usepackage[colorlinks=true,linkcolor=blue,citecolor=blue,urlcolor=blue]{hyperref}

\newtheorem{theorem}{Theorem}[section]
\newtheorem{lemma}[theorem]{Lemma}
\newtheorem{proposition}[theorem]{Proposition}
\newtheorem{corollary}[theorem]{Corollary}
\theoremstyle{definition}
\newtheorem{definition}[theorem]{Definition}
\newtheorem{hypothesis}[theorem]{Hypothesis}
\newtheorem{remark}[theorem]{Remark}

\newcommand{\lcm}{\operatorname{lcm}}
\newcommand{\sfree}{\operatorname{sf}}
\newcommand{\cR}{\mathcal R}
\newcommand{\cB}{\mathcal B}
\newcommand{\cQ}{\mathcal Q}
\newcommand{\cU}{\mathcal U}
\newcommand{\ckmin}{ck_{\min}}
\newcommand{\Lc}{\mathcal L}

% TODO(submission): author metadata.
\title[Large Erd\H{o}s--Straus witnesses]{Large Erd\H{o}s--Straus witnesses:\\
$\Omega$-results for the multiplier and slice parameters}
\author{Anonymous}
\date{Draft, 1 October 2026}

\begin{document}

\begin{abstract}
For a prime $p$ let $W(p)$ be the least modulus $M\equiv3\pmod4$ of a
``multiplier'' congruence that forces a solution of $4/p=1/x+1/y+1/z$, and
let $\ckmin(p)$ be the least product $ck$ of a positive Type~I slice.  Both
are expected to be at most a power of $\log p$ for almost all $p$, and are
known to exceed a constant times $\log p$ for infinitely many $p$.  We prove that
$W(p)\ge(\log p)^{2}\exp(-C\log\log p/\log\log\log p)$ for infinitely many
primes $p$ in Mordell's hard classes.  The construction is a prime-local
sieve: the class of one is imposed modulo all small primes, a congruence
saving bounds the residual sieve mass by $T^{o(1)}$, and the sieve is
transferred to primes by the Linnik-range prime number theorem of
Thorner and Zaman, whose relative error term absorbs a possible Siegel
zero.  We show that exponent~$2$ is the limit of every prime-local design
and isolate the exact combinatorial input needed to go further.  For the
slice parameter we prove that congruence methods certify exactly the least
quadratic non-residue, giving $\ckmin(p)\gg\log p\log\log\log p$
infinitely often (via Lau--Wu, after Graham--Ringrose) and showing that a
power saving would be a breakthrough on the least non-residue.  On the Haar
side we prove that the avoider density satisfies
$\log(1/\delta^*(T))\le T^{1/3+o(1)}$.
\end{abstract}

\maketitle

\section{Introduction}

The Erd\H{o}s--Straus conjecture asserts that $4/n=1/x+1/y+1/z$ is solvable
in positive integers for every $n\ge2$.  It suffices to treat primes, it has
been verified for $n\le10^{18}$ (\cite{PW}, \S6; \cite{Salez} for
$10^{17}$), and Mordell's identities settle every prime outside the six
classes $p\equiv1^2,11^2,13^2,17^2,19^2,23^2\pmod{840}$ \cite{Mordell}; we
call such primes \emph{hard}.  On the exceptional-set side, Vaughan
\cite{Vaughan} proved that the number $E(N)$ of $n\le N$ without a
representation satisfies $E(N)\ll N\exp\{-c(\log N)^{2/3}\}$; see
\cite{PW} for a uniform modern account and for many exceptions to the
analogous statement with large numerators, and \cite{ET} for the average
number of solutions, which over primes $p\le N$ is $(\log N)^3$ up to a factor
$\log\log N$.

This note is about the \emph{pointwise size} of the smallest witness.  Two
natural witness parameters are the following.

\smallskip\noindent\emph{The multiplier parameter.}  For $M\equiv3\pmod4$ put
$A_M=(M+1)/4$.  If $uvw=A_M$ and $n\equiv-uv^{-1}\pmod M$, then
$s=(nv+u)/M$ is a positive integer and
\begin{equation}\label{eq:ident}
 \frac4n=\frac1{suw}+\frac1{nsvw}+\frac1{nuvw}
\end{equation}
(direct verification; this is the classical multiplier family, cf.
\cite{Vaughan,ET}).  Define
\begin{equation}\label{eq:W}
 W(n)=\min\{M\equiv3\ (4):\ n\equiv-uv^{-1}\ (\mathrm{mod}\ M)\ \text{for some}\ uvw=A_M\},
\end{equation}
with $W(n)=+\infty$ if no such $M$ exists.  Thus $W(p)\le T$ supplies an
explicit representation with a modulus at most $T$.

\smallskip\noindent\emph{The slice parameter.}  A Type~I solution with
parameters $(a,b,c,k)$ satisfies $p(a+b)=k(4abc-1)$ and gives
$4/p=1/(ack)+1/(bck)+1/(pabc)$.  For an admissible pair $(c,k)$ (see
\S\ref{sec:typeI}) let $M_{c,k}(p)$ be the number of divisors of
$p^2+4ck^2$ lying in the class $-p\pmod{4ck}$; the slice is positive iff
$M_{c,k}(p)>0$.  Put $\ckmin(p)=\min\{ck:\ M_{c,k}(p)>0,\ \sfree(c)\notin\{1,2,3,6\}\}$
(the cores $1,2,3,6$ are excluded because they are deterministic on hard
primes).

\smallskip
Both parameters are expected to be small for almost all primes.  For $W$,
tails of the form
$\#\{p\le N:W(p)>T\}\ll N\exp\{-c(\log T)^2\log\log T\}$ (in the range
$\log N\gg(\log T)^3\log\log T$) are obtained in companion work whose
assembly is still under internal review; we do not use them.  In the other
direction, the class of one (every prime $p\equiv1\pmod{\lcm\{M\le T\}}$ has
$W(p)>T$) together with Linnik's theorem shows $W(p)\gg\log p$
infinitely often; with Chang's least-prime bound for smooth moduli
\cite{Chang} the constant becomes $5/8$.  The same is true for $\ckmin$ via
genus characters.  Whether $W(p)/\log p$ is unbounded was open.

\begin{theorem}[Proved modulo Theorem~\ref{thm:TZ}; effective]\label{thm:main}
There is an absolute constant $C$ such that for infinitely many hard primes $p$
\[
 W(p)\ \ge\ (\log p)^2\exp\Big(-C\,\frac{\log\log p}{\log\log\log p}\Big).
\]
More precisely, for every sufficiently large $T$ there is a prime
$p\equiv1\pmod{840}$ with $W(p)>T$ and
$\log p\le T^{1/2}\exp(O(\log T/\log\log T))$.
\end{theorem}

In particular the statement ``$W(p)\le(\log p)^A$ for all large $p$'' is
false for every $A<2$.  The cited input (Theorem~\ref{thm:TZ}) is Thorner and
Zaman's prime number theorem in the Linnik range \cite{TZ}, including the
McCurley zero-free region statement quoted there.  No hypothesis on Siegel
zeros is needed: the exceptional zero is absorbed by a uniqueness lemma
across the family of moduli (Lemma~\ref{lem:severe}) and a two-case
analysis (Theorem~\ref{thm:transfer}).

The exponent $2$ is the natural limit of the method
(Proposition~\ref{prop:ceiling}); beyond it one needs a pointwise
``hypergraph'' minorant, and Theorem~\ref{thm:Hmin} shows that such a
minorant would suffice, the prime side being completely handled.

For the slice parameter the picture is different: the only congruence
mechanism is the genus character, and it certifies exactly the least
quadratic non-residue $n_p$.

\begin{theorem}\label{thm:slice-intro}
\begin{enumerate}
\item[(a)] \textup{(Proved.)} $\ckmin(p)\ge n_p$ for every prime
$p\equiv1\pmod{24}$, and every fixed-period congruence minorant for
$\{\ckmin(p)>T\}$ is supported on primes with $n_p>T$
\textup{(Corollary~\ref{cor:exactly})}.
\item[(b)] \textup{(Proved modulo \cite[Prop.~5.1]{LW}.)}
$\ckmin(p)\gg\log p\cdot\log\log\log p$ for infinitely many hard primes.
\end{enumerate}
\end{theorem}

Thus a congruence proof of $\ckmin(p)>(\log p)^{1+\delta}$ infinitely often
would improve the best known unconditional (Graham--Ringrose) and
GRH-conditional (Montgomery) $\Omega$-results for $n_p$, and would
contradict the usual random model for $n_p$.

Finally, on the Haar (profinite) side, let $\delta^*(T)$ be the Haar measure
of the set of $n\in\widehat{\mathbb Z}^\times$ with $n\equiv1\pmod{24}$ and
$W(n)>T$, normalised within the class $1\pmod{24}$.  The class of one gives
$\log(1/\delta^*(T))\le(2/3+o(1))T$, and a crude prime-local bound gives
$T^{1/2+o(1)}$.

\begin{theorem}[Proved]\label{thm:haar-intro}
$\log(1/\delta^*(T))\le T^{1/3+o(1)}$.
\end{theorem}

Numerically $\log(1/\delta^*(T))$ grows like a small power of $\log T$, which
suggests that much larger values of $W(p)$ occur.  We do not pursue
heuristics here.

\smallskip\noindent\textbf{Status conventions.}  Every theorem header
carries a label.  \emph{Proved} means proved in full here.
\emph{Proved modulo X} means proved here assuming the cited statement X,
whose statement we checked in the source but whose proof we did not check.
\emph{Cited} marks external statements.

\section{The atoms of the multiplier problem}\label{sec:atoms}

For $M\equiv3\pmod4$ put $\cR(M)=\{-4D\bmod M:\ D\mid A_M^2\}$.

\begin{lemma}[Proved]\label{lem:atoms}
$\{-uv^{-1}\bmod M:\ uvw=A_M\}=\cR(M)$.  Consequently, for every integer $n$,
\[
 W(n)>T\iff n\bmod M\notin\cR(M)\ \text{for every}\ M\le T,\ M\equiv3\ (4).
\]
\end{lemma}

\begin{proof}
Since $4A_M\equiv1\pmod M$ and $\gcd(A_M,M)=1$, from $uvw=A_M$ we get
$v^{-1}\equiv4uw$, hence $-uv^{-1}\equiv-4u^2w$ and $D=u^2w$ divides $A_M^2$.
Conversely let $D\mid A_M^2$.  For each prime $q$ with $q^a\Vert A_M$ and
$q^d\Vert D$ ($d\le2a$) put $u_q=\max(0,d-a)$, $w_q=d-2u_q$, $v_q=a-u_q-w_q$;
these are non-negative, so $u=\prod q^{u_q}$, $v=\prod q^{v_q}$,
$w=\prod q^{w_q}$ satisfy $uvw=A_M$ and $u^2w=D$.
\end{proof}

\begin{lemma}[Class of one; proved]\label{lem:one}
$1\notin\cR(M)$ for every $M\equiv3\pmod4$.
\end{lemma}

\begin{proof}
If $-4D\equiv1$ then $D\equiv-A_M\pmod M$.  The map $D\mapsto A_M^2/D$
preserves this congruence, so we may assume $D\le A_M$; then
$0<D+A_M\le2A_M<M$ while $M\mid D+A_M$.
\end{proof}

\section{The prime-local residual system}\label{sec:local}

Fix $T$ large and put $\Lc=\log T$,
\begin{gather*}
 y=\sqrt T\,\exp(3\Lc/\log\Lc),\\
 Q=\lcm\big(24,\ \ell^{e_\ell}:\ \ell\le y\ \text{prime},\ \ell^{e_\ell}\le T\ \text{maximal}\big),
\end{gather*}
and let $\cU$ be the set of primes in $(y,T]$.  Then $\log Q\le2y$ for large
$T$ (only $\ell\le\sqrt T$ have $e_\ell\ge2$), and $840\mid Q$.  For
$\ell\in\cU$ let
\begin{gather*}
 F_\ell=\{-4D\bmod\ell:\ m\le T/\ell,\ m\ell\equiv3\ (4),\ D\mid A_{m\ell}^2,\
 m\mid 4D+1\},\\
 g(\ell)=\frac{|F_\ell|}{\ell-1},\qquad S=\sum_{\ell\in\cU}g(\ell).
\end{gather*}
All elements of $F_\ell$ are units.

\begin{lemma}[Prime-local reduction; proved]\label{lem:local}
If $n\equiv1\pmod Q$ and $n\bmod\ell\notin F_\ell$ for all $\ell\in\cU$,
then $W(n)>T$.
\end{lemma}

\begin{proof}
Let $M\le T$, $M\equiv3\pmod4$.  If every prime factor of $M$ is at most $y$,
then $M\mid Q$, so $n\equiv1\pmod M$ and Lemma~\ref{lem:one} applies.
Otherwise let $\ell>y$ divide $M$; since $\ell^2>T$ we have $\ell\Vert M$, and
$m=M/\ell<T/y\le y$, so $m\mid Q$ and $n\equiv1\pmod m$.  If
$n\equiv-4D\pmod M$ with $D\mid A_M^2$, then $m\mid4D+1$ and
$n\bmod\ell\in F_\ell$.  Now use Lemma~\ref{lem:atoms}.
\end{proof}

\begin{lemma}[Local sizes; proved]\label{lem:sizes}
$g(\ell)\le2T\tau^*(T)^2/y^2$, where $\tau^*(T)=\max_{n\le T}\tau(n)$.  In
particular $\max_\ell g(\ell)\le1/16$ for large $T$.
\end{lemma}

\begin{proof}
For each $m\le T/\ell$ there are at most $\tau(A^2)\le\tau(A)^2$ choices of
$D$, and $1/(\ell-1)\le2/\ell$.  Since $\log\tau^*(T)\le(\log2+o(1))\Lc/\log\Lc$
\cite[Thm.~317]{HW}, the bound is $\le2\exp((2\log2-6+o(1))\Lc/\log\Lc)$.
\end{proof}

The next lemma is the key arithmetic input: the congruence $m\mid4D+1$
saves a factor $m$.

\begin{lemma}[Mass bound; proved]\label{lem:mass}
With $X=(T+1)/4$,
\begin{align*}
 S&\le\tfrac43(3+\log X)\sum_{\substack{sr^2\le X\\ s\ \text{squarefree}}}\frac{\tau(4sr^2+1)}{sr}\\
 &\le 2(3+\Lc)(1+\Lc)^2\tau^*(T+2)=\exp\Big(O\Big(\frac{\log T}{\log\log T}\Big)\Big).
\end{align*}
\end{lemma}

\begin{proof}
Clearly $S\le\sum 1/(\ell-1)$ over triples $(m,\ell,D)$ as in the definition
of $F_\ell$; put $A=A_{m\ell}$.  Since $4A\equiv1$ and $4D\equiv-1\pmod m$, we
have $4A^2D^{-1}\equiv-1\pmod m$, so $D\mapsto A^2/D$ preserves $m\mid4D+1$
and exchanges $D>A$ with $D<A$; hence it suffices to count $D\le A$ and
double.  Write $D=sr^2$ with $s$ squarefree.  Then $D\mid A^2$ iff $sr\mid A$;
write $A=srk$, so that $D\le A$ iff $k\ge r$.  The map
$(m,\ell,D)\mapsto(m,s,r,k)$ is injective, and
$1/(\ell-1)\le2/\ell=2m/(4A-1)\le2m/(3A)$.

From $m\mid4sr^2+1$ and $m\mid m\ell=4srk-1$ we get $m\mid4sr(r+k)$, and
$\gcd(m,4sr)=1$, so
\[
 m\mid r+k .
\]
Since $k\ge r\ge1$, the least admissible $k$ is at least $\max(r,m-r)\ge m/2$,
so $\sum_{k\equiv-r\,(m),\ k\le X}1/k\le(3+\log X)/m$.  Summing over $m\mid4sr^2+1$
gives the first bound; the second follows from $\sum_{s\le X}1/s\cdot\sum_{r\le\sqrt X}1/r\le(1+\Lc)^2$
and $4sr^2+1\le T+2$.
\end{proof}

\begin{remark}
Without the congruence $m\mid r+k$ the same count gives $S\le T^{1/2-\eta+o(1)}$
for $y=T^{1/2+\eta}$, and the transfer below then only yields $W(p)\gg\log p$.
Numerically $S\approx 6.5,\ 13.6,\ 23.9,\ 38.5$ for $T=10^3,\dots,10^6$ (with
$y=\sqrt T$), i.e.\ $S$ is polylogarithmic in practice.
\end{remark}

\section{The analytic input}\label{sec:analytic}

\begin{theorem}[Thorner--Zaman {\cite[Cor.~1.4 and Rem.~1.5]{TZ}}; cited]\label{thm:TZ}
There are absolute effectively computable constants $c_4>0$, $C_0$ such
that for coprime $a,q$ with $q\ge2$ and $x\ge q^{12}$,
\begin{gather*}
 \sum_{\substack{p\le x\\ p\equiv a\,(q)}}\log p=\frac{\lambda x}{\varphi(q)}\,(1+\varepsilon),\\
 |\varepsilon|\le C_0\Big(e^{-c_4\log x/\log q}+e^{-c_4(\log x)^{3/5}/(\log\log x)^{1/5}}\Big),
\end{gather*}
where $\lambda=1-\chi_1(a)x^{\beta_1-1}/\beta_1$ if the exceptional zero $\beta_1$
for the modulus $q$ exists and $\lambda=1$ otherwise.  Here, following
McCurley \cite{McC} as quoted in \cite[p.~1]{TZ}, $\prod_{\chi\bmod q}L(s,\chi)$
has no zero in $\Re s\ge1-1/(13\log(q(|\Im s|+3)))$ apart from at most one
real zero $\beta_1$, which is simple and belongs to a unique real character
$\chi_1$; also $0<\lambda<2$.
\end{theorem}

We used the arXiv version; we did not compare the numbering with the
published version.  We also use the quoted McCurley statement for moduli
other than those of Theorem~\ref{thm:TZ}; alternatively the classical
Landau--Page theorem \cite[Ch.~14]{Dav} with a constant $c\le1/13$ in place
of $1/13$ works verbatim.

\begin{lemma}[One severe character per family; proved modulo the McCurley statement]\label{lem:severe}
Let $\cQ$ be a finite set of moduli, all $\le Z$.  Call $(\chi^*,\beta)$
\emph{severe} if $\chi^*$ is a real primitive character whose conductor $q^*$
divides some $q\in\cQ$, $L(\beta,\chi^*)=0$ and $\beta\ge1-1/(13\log(3Z^2))$.
Then there is at most one severe pair.  If $(\chi^*,\beta^*)$ is severe and
$q^*\mid q\in\cQ$, then $\beta_1=\beta^*$ and $\chi_1$ is induced by $\chi^*$ for
the modulus $q$.  If no severe pair has $q^*\mid q$, then any exceptional
$\beta_1$ of $q$ satisfies $x^{\beta_1-1}/\beta_1\le2\exp(-\log x/(13\log3Z^2))$.
\end{lemma}

\begin{proof}
If $(\chi_a,\beta_a),(\chi_b,\beta_b)$ are severe with conductors dividing
$q_a,q_b\in\cQ$, put $q'=\lcm(q_a,q_b)\le Z^2$.  Both characters induce
characters mod $q'$ with the same zeros in $\Re s>0$, and both $\beta$ lie in
McCurley's region for $q'$; uniqueness there gives $\beta_a=\beta_b$ and
$\chi_a=\chi_b$.  The second claim holds because $\beta^*$ lies in McCurley's
region for $q\le Z$.  For the third, an exceptional zero of $q$ in
$[1-1/(13\log 3Z^2),1)$ would be severe; also $\beta_1\ge1-1/(13\log6)>1/2$.
\end{proof}

\section{A transfer theorem}\label{sec:transfer}

\begin{theorem}[Proved modulo Theorem~\ref{thm:TZ}]\label{thm:transfer}
Let $T\ge1$, $Q\ge2$, and let $B(n)=\sum_{i\in I}c_i\,\mathbf 1[n\equiv b_i\ (d_i)]$
be a finite real combination with $\gcd(d_i,Q)=\gcd(b_i,d_i)=1$ and
$B(n)\le\mathbf 1[W(n)>T]$ for all $n\equiv1\pmod Q$.  Put
$Z=Q\max d_i$, $\mu=\sum c_i/\varphi(d_i)$, $M_1=\sum|c_i|/\varphi(d_i)$, and assume
$\mu>0$ and the \emph{twist condition}: for every real primitive character
$\psi$ of conductor $f>1$ with $\gcd(f,Q)=1$ and $f\mid d_i$ for some $i$,
\[
 |\mu_\psi|\le\mu/4,\qquad \mu_\psi=\sum_{i:\,f\mid d_i}c_i\psi(b_i)/\varphi(d_i).
\]
Let $K=1+\log(M_1/\mu)$.  There is an absolute $C_1$ such that if
$\log x\ge C_1K\max(\log Z,K)$, then some prime $p\le x$, $p\equiv1\pmod Q$,
has $W(p)>T$.
\end{theorem}

\begin{proof}
Let $q_i=Qd_i$ and let $a_i$ be the class $\equiv1\ (Q)$, $\equiv b_i\ (d_i)$.
Then $\sum_{p\le x,\,p\equiv1(Q)}\log p\,B(p)=\sum_ic_i\theta(x;q_i,a_i)$, and
$x\ge Z^{12}\ge q_i^{12}$.  Put
$E_1=e^{-c_4\log x/\log Z}+e^{-c_4(\log x)^{3/5}/(\log\log x)^{1/5}}$ and
$E_2=e^{-\log x/(13\log3Z^2)}$, and apply Lemma~\ref{lem:severe} to
$\{q_i\}$.  Write the conductor of a severe $\chi^*$ as $q^*=q^*_Qq^{*\prime}$ with
$q^*_Q$ supported on the primes of $Q$ and $\gcd(q^{*\prime},Q)=1$, and
$\chi^*=\chi^*_Q\chi^{*\prime}$.

\emph{Case A: a severe pair exists with $q^*\mid Q$.}  Then for every $i$,
$\lambda_i=1-\chi^*(a_i)x^{\beta^*-1}/\beta^*=1-x^{\beta^*-1}/\beta^*=:\lambda_*>0$,
because $a_i\equiv1\pmod{q^*}$; so
$\sum c_i\theta\ge(x\lambda_*/\varphi(Q))(\mu-C_0E_1M_1)$.

\emph{Case B: otherwise.}  If $q^*\mid q_i$ (equivalently $q^{*\prime}\mid d_i$,
$q^{*\prime}>1$), then $\lambda_i=1-\psi(b_i)x^{\beta^*-1}/\beta^*$ with
$\psi=\chi^{*\prime}$; otherwise $|\lambda_i-1|\le2E_2$.  Since $\lambda_i<2$ and
$x^{\beta^*-1}/\beta^*\le2$,
\begin{align*}
 \sum c_i\theta&\ge\frac{x}{\varphi(Q)}\big(\mu-2|\mu_\psi|-M_1(2C_0E_1+2E_2)\big)\\
 &\ge\frac{x}{\varphi(Q)}\big(\mu/2-M_1(2C_0E_1+2E_2)\big).
\end{align*}
Both brackets are positive once $E_1,E_2\le\mu/(8(C_0+1)M_1)$, which holds
under the stated condition (use $u^{3/5}/(\log u)^{1/5}\ge u^{1/2}$).  Hence
$\sum_p\log p\,\mathbf 1[W(p)>T]\ge\sum_p\log p\,B(p)>0$.
\end{proof}

No lower bound for $1-\beta^*$ is used: in Case~A the exceptional factor
$\lambda_*$ multiplies main term and error alike.  A prime number theorem
valid only for $q\le\exp(c\sqrt{\log x})$ would force $\log x\gg(\log Z)^2\ge T$;
the Linnik-range uniformity is essential.

\section{Proof of Theorem~\ref{thm:main}}\label{sec:proof}

Let $J$ be the least even integer $\ge22S+10$ and
\begin{gather*}
 B(n)=\sum_{P\subseteq\cU,\ |P|\le J-1}(-1)^{|P|}\prod_{\ell\in P}\mathbf 1[n\bmod\ell\in F_\ell]
 =\sum_{j=0}^{J-1}(-1)^j\binom{N(n)}j,\\
 N(n)=\#\{\ell\in\cU:\ n\bmod\ell\in F_\ell\}.
\end{gather*}
\emph{Pointwise bound.}  For $N\ge1$,
$\sum_{j\le J-1}(-1)^j\binom Nj=(-1)^{J-1}\binom{N-1}{J-1}\le0$; hence
$B\le\mathbf 1[N=0]\le\mathbf 1[W>T]$ on $n\equiv1\ (Q)$ by Lemma~\ref{lem:local}.
Expanding $\mathbf 1[n\bmod\ell\in F_\ell]=\sum_{a\in F_\ell}\mathbf 1[n\equiv a\ (\ell)]$ puts $B$
in the form of Theorem~\ref{thm:transfer}, with squarefree moduli
$d_P=\prod_{\ell\in P}\ell\le T^{J}$ coprime to $Q$.

\emph{Mean and mass.}  With $e_j$ the elementary symmetric functions of the
$g(\ell)$, $\mu=\sum_{j<J}(-1)^je_j$ and $M_1=\sum_{j<J}e_j\le e^S$.  Let $N$ be a sum
of independent Bernoulli$(g(\ell))$ variables; the identity above and
$\binom{N-1}{J-1}\le\binom NJ$ give
$|\mu-V|\le e_J\le(eS/J)^J\le e^{-2J}$, where $V=\prod(1-g(\ell))\ge e^{-2S}$
(Lemma~\ref{lem:sizes}).  So $\mu\ge0.99V$ and $\log(M_1/\mu)\le3S+1$.

\emph{Twist.}  If $\psi$ is real primitive of conductor $f>1$, coprime to $Q$,
dividing some $d_P$, then $f=\prod_{\ell\in P_f}\ell$ with $r=|P_f|\ge1$ and
$\psi=(\cdot/f)$.  With $h(\ell)=\sum_{a\in F_\ell}(a/\ell)/(\ell-1)$, $|h|\le g$,
\[
 \mu_\psi=(-1)^r\prod_{\ell\in P_f}h(\ell)\sum_{P'\subseteq\cU\setminus P_f,\ |P'|\le J-1-r}(-1)^{|P'|}\prod_{P'}g ,
\]
which vanishes if $r\ge J$.  The inner sum is $V'+\varepsilon'$ with
$V'\le(16/15)^rV$ and $|\varepsilon'|\le S^{J-r}/(J-r)!$, so
$|\mu_\psi|\le V/15+e^{-15S}\le\mu/4$ (treat $r\le J/2$ and $r>J/2$ separately,
using $g\le1/16$ and $(e/11)^{1/2}<e^{-0.69}$).

\emph{Moduli.}  $\log Z\le\log Q+J\Lc\le3y$ because $S=T^{o(1)}$
(Lemma~\ref{lem:mass}).  Theorem~\ref{thm:transfer} gives a prime $p\equiv1\pmod Q$
with $W(p)>T$ and $\log p\le C_1(3S+2)\cdot3y\le\sqrt T\exp((3+\log2+o(1))\Lc/\log\Lc)$.
Since $840\mid Q$ the prime is hard, and $p>Q>T$ gives infinitely many distinct
primes.  Inverting, $\Lc\ge2\log\log p-O(\log\log p/\log\log\log p)$.  All
inputs are effective.\qed

\section{Limits of the method}\label{sec:limits}

\subsection{Complete certificates}

A \emph{complete certificate} for $W>T$ is a reduced class $c\bmod Q$ all of
whose primes have $W(p)>T$.  The class of one modulo
$L^*(T)=\lcm(24,\ M\le T,\ M\equiv3\ (4))$ is one, with
$\log L^*(T)=(2/3+o(1))T$.

\begin{proposition}[Proved]\label{prop:complete}
Every complete certificate has $\log Q\ge(2/3-o(1))T$.
\end{proposition}

\begin{proof}
(i) If a prime $\ell\equiv3\pmod4$, $\ell\le T$, does not divide $Q$, then $c$
is compatible with the atom $-4\in\cR(\ell)$ ($D=1$), and Dirichlet's theorem
gives primes with $W\le\ell$; so $\prod_{\ell\le T,\ell\equiv3(4)}\ell\mid Q$.
(ii) Since $\cR(3)=\{2\}$, $3\mid Q$ and $c\equiv1\pmod3$.  (iii) Let
$\ell\equiv1\pmod4$, $\ell\le T/3$, and suppose a prime $r\equiv2\pmod3$ divides
$A=(3\ell+1)/4$.  The atom $M=3\ell$, $D=r$ has class $-4r\equiv1\pmod3$; if
$\ell\nmid Q$ it is compatible with $c$, a contradiction.  (iv) The primes
$\ell\equiv1\pmod4$ for which $(3\ell+1)/4$ has no prime factor $\equiv2\pmod3$
avoid the classes $\ell\equiv-3^{-1}\pmod r$ for all primes $r\equiv2\pmod 3$,
$r\ge5$; by the prime number theorem in progressions their logarithmic
weight up to $X$ is $o(X)$.  Hence $\log Q\ge T/2+T/6-o(T)$.
\end{proof}

With Chang's theorem this gives the previous record $W\ge(5/8-\varepsilon)\log p$,
and no complete certificate can do better than linear.

\subsection{Prime-local designs}

Call a design \emph{prime-local} if it imposes $n\equiv1\pmod{\ell^{e_\ell}}$ for
$\ell$ in a set $\Pi$ and every $M\le T$, $M\equiv3\pmod4$, has at most one prime
factor outside $\Pi$.  Theorem~\ref{thm:main} uses such a design.

\begin{proposition}[Proved]\label{prop:ceiling}
In a prime-local design, $\Pi$ contains every odd prime $\le\sqrt{T/3}$ with at
most one exception.  Hence every prime it produces has
$\log p\ge(1/\sqrt3-o(1))\sqrt T$.
\end{proposition}

\begin{proof}
For odd primes $\ell_1<\ell_2\le\sqrt{T/3}$ outside $\Pi$, one of
$\ell_1\ell_2,3\ell_1\ell_2$ is $\equiv3\pmod4$ and $\le T$.
\end{proof}

\subsection{What would suffice beyond exponent two}

\begin{hypothesis}[$H_{\min}(\theta)$; open for $\theta<1/2$]\label{hyp:min}
For every $\varepsilon>0$ and large $T$ there are $Q$ with $840\mid Q$,
$\log Q\le T^{\theta+\varepsilon}$, and a minorant $B$ as in
Theorem~\ref{thm:transfer} with $\log\max d_i\le T^{\theta+\varepsilon}$,
$\log(M_1/\mu)\le T^{\varepsilon}$, $\mu>0$, and the twist condition.
\end{hypothesis}

\begin{theorem}[Proved modulo Theorem~\ref{thm:TZ}]\label{thm:Hmin}
$H_{\min}(\theta)$ implies $W(p)>(\log p)^{1/\theta-\varepsilon}$ for infinitely
many hard primes, for every $\varepsilon>0$.
\end{theorem}

\begin{proof}
Adjoin to $Q$ a prime $\ell_0\in(R,2R]$, $R=\max(T,\max d_i)$; this keeps all
hypotheses and forces $p>T$.  Theorem~\ref{thm:transfer} gives
$\log p\le T^{\theta+2\varepsilon+o(1)}$.
\end{proof}

$H_{\min}$ is a purely combinatorial statement about congruences.  For
$\theta<1/2$ the residual atoms involve several free primes, and plain
event-level Bonferroni truncation over the raw atom list fails: configurations in which $n\equiv-4$
modulo many free primes make $\binom{N-1}{J-1}$ large on a set of measure
$\ge T^{-O(\sqrt J)}$, so the truncated mean is negative for
$4.2\theta^2(\log T)^2\lesssim J\le T^\theta$.  A hub-adapted (hypergraph)
minorant would be needed.

\section{The slice parameter}\label{sec:typeI}

For a prime $p$ let $\cB_p$ be the set of pairs $(c,k)$ with
$1\le k\le\lfloor2p/3\rfloor$, $1\le c\le\lfloor(2p+k)/(4k)\rfloor$, $(p,ck)=1$;
put $h=4ck$, $N_{c,k}=p^2+4ck^2$, $s=\sfree(c)$, and let $\chi_s=(\Delta_s/\cdot)$,
where $\Delta_s<0$ is the discriminant of $\mathbb Q(\sqrt{-s})$; its
conductor divides $4s\mid h$.  $M_{c,k}(p)$ is the number of divisors of
$N_{c,k}$ in the class $-p\pmod h$; a divisor $e$ in that class gives the
Type~I solution $a=(p+e)/h$, $b=(p+N/e)/h$.

\begin{lemma}[Genus obstruction; proved]\label{lem:genus}
If $p\equiv1\pmod{24}$, $(c,k)\in\cB_p$ and $\chi_s(p)=1$, then $M_{c,k}(p)=0$.
\end{lemma}

\begin{proof}
Write $c=st^2$.  Every prime $q\mid N$ is odd, prime to $stk$, and
$(p(2tk)^{-1})^2\equiv-s\pmod q$, so $\chi_s(q)=1$; hence $\chi_s(e)=1$ for every
divisor $e$ of $N$.  But $\chi_s(-p)=\chi_s(-1)\chi_s(p)=-1$, and $\chi_s$ is
defined modulo $h$.
\end{proof}

\begin{lemma}[A sufficient prime factor; proved]\label{lem:goodprime}
If a prime $q$ divides $N_{c,k}$ and $q\equiv-p\pmod h$, then $M_{c,k}(p)\ge1$.
\end{lemma}

\begin{proof}
$q$ itself is a divisor in the class $-p$.
\end{proof}

\begin{lemma}[Proved]\label{lem:np}
Let $p\equiv1\pmod{24}$ and suppose $(\ell/p)=1$ for every prime $5\le\ell\le B$.
Then $\ckmin(p)>B$.  In particular $\ckmin(p)\ge n_p$, the least quadratic
non-residue modulo $p$.
\end{lemma}

\begin{proof}
For $ck\le B$, $s\le B$ and
$\chi_s(p)=(-s/p)=(-1/p)\prod_{q\mid s}(q/p)=1$, since
$(-1/p)=(2/p)=(3/p)=1$.  Apply Lemma~\ref{lem:genus}.
\end{proof}

\begin{proposition}[Proved]\label{prop:QR}
Let $24\mid L$ and let $a\bmod L$ be a reduced class with $a\equiv1\pmod{24}$
such that every sufficiently large prime $p\equiv a\pmod L$ has $\ckmin(p)>T$.
Then every prime $5\le\ell\le T$ divides $L$ and $(a/\ell)=1$.
\end{proposition}

\begin{proof}
Suppose $\ell\nmid L$ or $(a/\ell)=-1$.  Choose $c_0\bmod4\ell$ with
$c_0\equiv a\pmod{\gcd(L,4\ell)}$, $c_0\equiv1\pmod4$, $(c_0/\ell)=-1$; for
$n\equiv1\pmod 4$, $\chi_\ell(n)=(n/\ell)$, so $\chi_\ell(c_0)=-1$.  Take a prime
$q\equiv-c_0\pmod{4\ell}$, $q\nmid L$.  Then $\chi_\ell(q)=\chi_\ell(-1)\chi_\ell(c_0)=1$,
so $-4\ell\equiv\rho^2\pmod q$ for some $\rho\not\equiv0$.  The conditions
$p\equiv a\ (L)$, $p\equiv c_0\ (4\ell)$, $p\equiv\rho\ (q)$ are compatible and
define a reduced class, which contains infinitely many primes.  For these,
$q\mid p^2+4\ell$ and $q\equiv-p\pmod{4\ell}$, so $M_{\ell,1}(p)\ge1$
(Lemma~\ref{lem:goodprime}), and $(\ell,1)\in\cB_p$ for large $p$, whence
$\ckmin(p)\le\ell\le T$.
\end{proof}

\begin{corollary}[Proved]\label{cor:exactly}
Let $B(n)$ be a finite combination of residue classes with
$B(p)\le\mathbf 1[\ckmin(p)>T]$ for all large primes $p$ in a reduced class
$\equiv1\pmod{24}$.  Then every large prime $p$ in that class with $B(p)>0$ has
$n_p>T$.  Together with Lemma~\ref{lem:np}: the depth certifiable by
congruences is exactly $n_p$.
\end{corollary}

\begin{proof}
$B$ is periodic modulo some $L$; if $B(p)>0$ then $B$ is positive on the
whole class $p\bmod L$, so Proposition~\ref{prop:QR} gives $(\ell/p)=1$ for
$5\le\ell\le T$, and $(2/p)=(3/p)=1$.
\end{proof}

\begin{theorem}[Proved modulo {\cite[Prop.~5.1]{LW}}]\label{thm:slice}
$\ckmin(p)\gg\log p\cdot\log\log\log p$ for infinitely many hard primes $p$.
\end{theorem}

\begin{proof}
Lau and Wu \cite[Prop.~5.1]{LW} (following Graham--Ringrose \cite{GR}) show:
for small fixed $\delta>0$ and $y(x)=\delta\log x\log_3x$ there are
$x_n\to\infty$ and primes $p\in(x_n^{1/2},x_n\log x_n]$ with $p\equiv1\pmod4$
and $(q/p)=1$ for all primes $q\le y(x_n)$.  Such $p$ are $\equiv1\pmod{24}$
and, since $(5/p)=(7/p)=1$, lie in a Mordell class.  Lemma~\ref{lem:np} gives
$\ckmin(p)>y(x_n)\ge(\delta/2)\log p\log_3p$.  Effectivity is not claimed.
\end{proof}

Under GRH, Montgomery's $n_p=\Omega(\log p\log\log p)$ would give the
corresponding bound, and Ankeny's $n_p\ll(\log p)^2$ caps what congruences can
certify; the random model predicts $\max_{p\le x}n_p\asymp\log x\log\log x$
(statements as quoted in \cite{LW}).  Finally, the primes of
Theorem~\ref{thm:main} are quadratic residues of every $\ell\le y$, so
Lemma~\ref{lem:np} gives, \emph{jointly}, $W(p)\ge(\log p)^{2-o(1)}$ and
$\ckmin(p)\ge(\log p)^{1-o(1)}$ for infinitely many hard $p$ (proved modulo
Theorem~\ref{thm:TZ}).

\section{The Haar side}\label{sec:haar}

Impose the class of one at all primes $\le z$ ($z\ge3$), i.e.\ $n\equiv1\pmod{Q_z}$
with $Q_z$ defined like $Q$.  For $M\le T$ write $M=mr$ with $m$ the $z$-smooth
and $r$ the $z$-rough part.  If $r>1$, the surviving atoms of $M$ are the
$D\mid A_M^2$ with $m\mid4D+1$, each giving the event $n\equiv-4D\pmod r$ of
conditional Haar probability $1/\varphi(r)$.

\begin{lemma}[Proved]\label{lem:struct}
If $M\equiv3\pmod4$, $D\mid A_M^2$ and $M=mr$ with $m\mid4D+1$, then $m\le r^2+1$.
\end{lemma}

\begin{proof}
By the involution of Lemma~\ref{lem:mass} assume $D\le A$, $D=sr'^2$,
$A=sr'k$, $k\ge r'$.  As there, $m\mid r'+k\le2k$, so
$r=(4sr'k-1)/m\ge2sr'-1/(2k)$, i.e.\ $r\ge2sr'$, and
$m\le4sr'^2+1\le r^2+1$.
\end{proof}

\begin{lemma}[Proved]\label{lem:globalmass}
For every $3\le z\le T$, the total mass $S_{\rm tot}=\sum1/\varphi(r)$ over all
surviving atoms is $\exp(O(\log T/\log\log T))$.
\end{lemma}

\begin{proof}
Use $1/\varphi(r)\ll\log\log T\cdot m/M$ and repeat the proof of
Lemma~\ref{lem:mass} with $\ell$ replaced by $r$; every surviving atom has
$m\mid\gcd(4D+1,M)$, hence $m\mid r'+k$.
\end{proof}

(Using \cite[Prop.~1.4]{ET} on dyadic blocks one gets
$S_{\rm tot}\ll(\log T)^4\log\log T$; this is not needed below.)

\begin{proof}[Proof of Theorem~\ref{thm:haar-intro}]
Fix $\varepsilon>0$, put $y=T^{1/3+\varepsilon}$ and $z=y$.  Every rough part is
$\ell$, $\ell^2$ or $\ell_1\ell_2$.  Use the independent coordinates
$X_\ell=n\bmod\ell^{e_\ell}$.  Single atoms form forbidden sets $G_\ell$ of
measure $g_\ell$, with $1\notin G_\ell$ (Lemma~\ref{lem:one}) and
$\sum g_\ell\le S_{\rm tot}$.

\emph{Bad primes.}  $B=\{\ell:\ g_\ell>1/4\}$ has $|B|\le4S_{\rm tot}=T^{o(1)}$.
Impose also $n\equiv1\pmod{\ell^{e_\ell}}$ for $\ell\in B$ (cost $\le|B|\log T$).
Pairs inside $B$ die by Lemma~\ref{lem:one}; a pair atom $(\ell b,a)$ with
$b\in B$ becomes the residue $a\bmod\ell$ at $\ell$ if $a\equiv1\ (b)$ and dies
otherwise.  The enlarged sets satisfy
$g'_\ell\le1/4+|B|\tau^2_{\max}T/(\ell y(\ell-1))\le1/2$, with
$\tau^2_{\max}=\max_{A\le T}\tau(A^2)=T^{o(1)}$, and
$\sum(g'_\ell-g_\ell)\le|B|\tau^2_{\max}T/y^2=T^{1/3-2\varepsilon+o(1)}$.

\emph{Good pairs.}  Under the product measure $\mu'$ uniform on the allowed
residues, a good pair event has $\mu'(E)\le4/((\ell_1-1)(\ell_2-1))$, and the
per-prime weight is
$w'_\ell\le16\tau^2_{\max}T/(\ell^2y)\le T^{-3\varepsilon+o(1)}$.  The asymmetric local
lemma \cite[Lemma~5.1.1]{AS} with $x_E=2\mu'(E)$ gives
$\mu'(\text{no good pair event})\ge\exp(-16S_{\rm tot})$.

\emph{Collecting.}  With $Q'=Q_y\prod_{b\in B}b^{e_b}$,
\begin{align*}
\delta^*(T)&\ge\frac{8}{\varphi(Q')}\prod(1-g'_\ell)\,e^{-16S_{\rm tot}}\\
&\ge\exp\big(-\pi(y)\log T-|B|\log T-18S_{\rm tot}-T^{1/3-2\varepsilon+o(1)}\big)\\
&=\exp\big(-T^{1/3+\varepsilon+o(1)}\big).
\end{align*}
\end{proof}

The exponent $1/3$ reflects the crude bound $N(M)\le\tau(A^2)$ for multi-prime
atoms at a fixed prime.  A congruence saving \emph{uniform in a fixed prime}
(the per-prime analogue of Lemma~\ref{lem:mass}) would, by the local lemma,
give $\log(1/\delta^*(T))\le(\log T)^{O(1)}$; we do not know how to prove it.
None of this transfers to primes: the local lemma does not provide a
pointwise minorant.

\section{Summary of status}

\begin{itemize}
\item Proved: Lemmas~\ref{lem:atoms}--\ref{lem:mass}, \ref{lem:genus}--\ref{lem:np},
\ref{lem:struct}, \ref{lem:globalmass}; Propositions~\ref{prop:complete},
\ref{prop:ceiling}, \ref{prop:QR}; Corollary~\ref{cor:exactly};
Theorem~\ref{thm:haar-intro} (with the classical local lemma and divisor bound).
\item Proved modulo Thorner--Zaman \cite{TZ} (with the McCurley statement quoted there):
Lemma~\ref{lem:severe}, Theorems~\ref{thm:transfer}, \ref{thm:main}, \ref{thm:Hmin}, and the joint
$W$/$\ckmin$ statement.
\item Proved modulo Lau--Wu \cite[Prop.~5.1]{LW}: Theorem~\ref{thm:slice}.
\item Open: $H_{\min}(\theta)$ for $\theta<1/2$; a polylogarithmic Haar bound;
any bound $W(p)>(\log p)^{2+\delta}$ or $\ckmin(p)>(\log p)^{1+\delta}$
infinitely often.
\end{itemize}
Nothing here bears on the truth of the Erd\H{o}s--Straus conjecture: a prime with
large $W(p)$ or $\ckmin(p)$ may well have other representations.

\begin{thebibliography}{99}
\bibitem{AS} N.~Alon and J.~H.~Spencer, The Probabilistic Method, 4th ed.,
Wiley, 2016.
\bibitem{Chang} M.-C.~Chang, Short character sums for composite moduli,
J.\ Anal.\ Math.\ 123 (2014), 1--33.
\bibitem{Dav} H.~Davenport, Multiplicative Number Theory, 3rd ed., Springer GTM 74, 2000.
\bibitem{ET} C.~Elsholtz and T.~Tao, Counting the number of solutions to the
Erd\H{o}s--Straus equation on unit fractions, J.\ Aust.\ Math.\ Soc.\ 94 (2013),
50--105 (Proposition numbering as in arXiv:1107.1010v6).
\bibitem{GR} S.~W.~Graham and C.~J.~Ringrose, Lower bounds for least quadratic
non-residues, in: Analytic Number Theory (Allerton Park, 1989), Progr.\ Math.\ 85,
Birkh\"auser, 1990, 269--309.
\bibitem{HW} G.~H.~Hardy and E.~M.~Wright, An Introduction to the Theory of
Numbers, 6th ed., Oxford, 2008.
\bibitem{LW} Y.-K.~Lau and J.~Wu, On the least quadratic non-residue,
Int.\ J.\ Number Theory 4 (2008), 423--435.
\bibitem{McC} K.~S.~McCurley, Explicit zero-free regions for Dirichlet
$L$-functions, J.\ Number Theory 19 (1984), 7--32.
\bibitem{Mordell} L.~J.~Mordell, Diophantine Equations, Academic Press, 1969,
Chapter~30.
\bibitem{PW} C.~Pomerance and A.~Weingartner, Exceptions to the
Erd\H{o}s--Straus--Schinzel conjecture, arXiv:2511.16817 (2025).
\bibitem{Salez} S.~E.~Salez, The Erd\H{o}s--Straus conjecture: new modular
equations and checking up to $N=10^{17}$, arXiv:1406.6307 (2014).
\bibitem{TZ} J.~Thorner and A.~Zaman, Refinements to the prime number theorem
for arithmetic progressions, Math.\ Z.\ 306 (2024), Paper No.~54
(arXiv:2108.10878v2).
\bibitem{Vaughan} R.~C.~Vaughan, On a problem of Erd\H{o}s, Straus and
Schinzel, Mathematika 17 (1970), 193--198.
\end{thebibliography}

\end{document}
